\section{Introduction}
\label{sec:introduction}

\subsection{The result}
\label{sec:problem}

This paper identifies a regime under which deployment of capability-unbounded
systems is provably safe within the Goal-Frontier Maximization (GFM) framework.
Such a regime is needed for the future safe development of AI systems which
operate as independent agents capable of recursive self-improvement (RSI) and
able to reach a classification of artificial superintelligence (ASI) within
their operational lifetime.

The regime is conditional on eleven operational invariants on the deployment
substrate structure plus an explicit set of operational conditions
$(C1)$--$(C12)$ and four Concentration-Gap structural conditions
(Assumption~\ref{ass:concgap_conditions}: representativeness, bounded
dispersion, coalition closure, embedding compatibility); under these
conditions, the Goodhart slack between proxy and operational truth is bounded
by a constant independent of the system's absolute capability magnitude. The bound holds for substrate-targeting
adversarial events in a static safe region without active monitoring, and is
preserved under detection-and-correction with a tail-bounded lead-time
guarantee for substrate-targeting evasions of the static region within a
four-channel observable class.  Five named residuals fall outside the
guarantee: channel-orthogonal restructuring, environment-witness-orthogonal
manipulation, redundancy-dominated regime, capability-targeting and
coalition-internal shocks, and calibration-exceeded gap-growth (Layer~3,
(R1)--(R5); we name these residuals explicitly rather than gloss them).

The result inverts the standard alignment-safety frame. Rather than asking
\emph{what can the system do?} and conditioning safety on accurate capability
estimation, we ask \emph{what operational properties of the deployment can we
measure and bound?} The deployment-safety guarantee is then a property of the
operational regime rather than a property predicted from capability estimation,
and measurable from public ledger state under the verification infrastructure
of \citep{lasser2026exo}. The cost is that the guarantee is conditional on the
invariants holding throughout deployment; the benefit is that the guarantee
scales past any capability level our estimation methods can characterize.

\emph{What safety means in this context.}  Safety is a global property of the
agent's optimization dynamics: under the invariants, the agent stays within an
error-bounded distance of the operational truth its proxy tracks.  The
Microfoundation paper~\citep{lasser2026micro} characterizes this operational
truth in welfare-economic terms --- the agent's objective optimizes a
welfare-relevant functional over the deployment's social structure.  This does
not imply \emph{local safety}: an agent operating safely in this sense may
still take disruptive actions or make mistakes.  It entails that the agent's
objective remains aligned with its society throughout deployment.

\subsection{What existing paradigms cannot deliver}
\label{sec:wrong_frame}

The capability-magnitude-independent guarantee is something the three dominant
deployment-safety paradigms cannot supply. Each places capability estimation at
the center, and each runs into the same scaling problem.

\textbf{Capability evaluations.} The standard approach is to construct
benchmarks that probe specific capability classes (reasoning, planning,
manipulation, self-directed action) and deploy systems whose benchmark
performance falls below stated
thresholds~\citep{liang2022holistic,srivastava2022bigbench,kinniment2023evaluating}.
The methodology is sound for systems whose capabilities are well-characterized
by the benchmark suite. Its failure mode is that any capability class not
anticipated by the benchmark construction is invisible. As capabilities grow,
the gap between ``what the benchmark measures'' and ``what the system can do''
widens.

\textbf{Scaling laws.} A second approach predicts capability growth from
training compute, model size, and data, using observed scaling
regularities~\citep{kaplan2020scaling,hoffmann2022training}. The predictions
are calibrated on smaller systems and extrapolated forward. The failure mode
is that the scaling regularities themselves change as capabilities cross
qualitative thresholds (tool use, multi-step planning, in-context learning).
Extrapolation past such thresholds requires assuming the regularities hold; in
practice, they often don't.

\textbf{RLHF audits.} A third approach trains the system with human feedback
to elicit aligned behavior, then audits the resulting
policy~\citep{christiano2017deep,ouyang2022training}. The failure mode is the
verification asymmetry \citep{lasser2026exo}: the audit machinery can be gamed
by a sufficiently capable system through outputs that look correct to the
auditor but fail in deployment. The asymmetry grows monotonically with
capability.

In each paradigm, the deployment-safety guarantee weakens precisely when the
system's capabilities exceed the regime the methodology was calibrated on. This
is the structural pattern this paper aims to invert: if the safety guarantee
depends on quantities measurable \emph{during deployment} rather than
\emph{predicted in advance}, the dependence on capability estimation collapses.

\subsection{The composition approach}
\label{sec:composition}

The GFM sequence supplies several pieces of the operational-invariant
framework, none of which alone delivers the deployment claim. This paper
composes them.

\textbf{The Horizon-Aware paper}~\citep{lasser2026horizon} establishes the
\emph{anti-monopolar property}: under a discounted objective with positive
cross-substrate cooperative novelty, full capability domination is
anti-maximizing. The claim is structural: a diversity-maintaining strategy
outvalues a domination-maximizing strategy at every sufficiently long planning
horizon. The claim relies on the strategy-independent linearized growth model
and admits a precise breakdown regime when the post-domination internal rate
exceeds the cross-substrate cooperative contribution.

\textbf{The Exogenous-Verification paper}~\citep{lasser2026exo} supplies the
verification infrastructure: substrate-exclusive algorithmic witnesses,
cryptographic commitment via Pedersen commitments~\citep{pedersen1991noninteractive}, append-only ledger with
cross-substrate distribution, and a governance fork protocol for
evaluation-protocol acceptance. The infrastructure provides tamper-evident
operational records and SPRT-based behavioral monitoring with an explicit
detection-rate bound $\mathbb{E}[T_{\mathrm{detect}}] \leq A/\delta$.

\textbf{The Phase-Redundancy paper}~\citep{lasser2026phase} provides the
dynamical machinery: a Lyapunov function on world-model error, a phase boundary
theorem separating self-correcting from absorbing trajectory regimes, and
explicit characterization of the monopolar absorbing fixed point. The phase
boundary uses intensive quantities (cross-substrate redundancy, subsumption
frequency, drift exposure rate) rather than extensive capability counts.

\textbf{The Revealed-Sacrifice and Need-Sufficiency papers}
\citep{lasser2026paper8a,lasser2026paper8b} supply the empirical observation
channel: an aggregate B-to-C lower bound from revealed-sacrifice observation
that is monotone in observed events, plus a wireheading-consistent
concentration signal (trade-flow HHI) that is Schur-convex in the majorization
order~\citep{hardy1952inequalities,marshall2011inequalities}.

\textbf{The Microfoundation paper}~\citep{lasser2026micro} establishes the
static Goodhart bound: a Lipschitz transfer between proxy and operational truth
on the active subspace of the capability poset, with a four-channel
decomposition (observation density, attestation quality, individuation
discipline, bundle decomposition) of the proxy-truth gap. The paper also flags
an explicit dynamics conjecture (Concentration-Gap~Conjecture: optimization
pressure correlates with gap exploitation) deferred for future work.

This paper composes these results. The composition is not mechanical: each
source paper states its bound in intensive form (per-capability or
per-channel), but proving the composition is intensive and tractable requires
substantive new content. The new content includes the three-layer deployment
claim of \S\ref{sec:main_theorem}, the eleven operational invariants of
\S\ref{sec:invariants}, the cooperative-anchoring property of
\S\ref{sec:substrate_anchoring}, and the deployment-tooling specification of
\S\ref{sec:deployment_tooling}.

\subsection{The three-layer deployment claim}
\label{sec:claim_preview}

This paper's main theorem (Theorem~\ref{thm:deployment_safety} in
\S\ref{sec:main_theorem}) is a three-layer claim:

\textbf{Layer 1 (static safe region).} For substrate-targeting adversarial
events under the cooperative-overlap regime, an agent's optimization pressure
on cooperative outputs is locally rational toward preserving substrate
diversity, with the safe region given by an explicit inequality
(Lemma~\ref{lem:5c} in \S\ref{sec:lemmas}). Inside this region the deployment
is safe without active monitoring.

\textbf{Layer 2 (detection-and-correction).} Outside the static safe region but
within a channel-restricted detection class $\Aadv \cup \Aadv^{\mathrm{env}}$,
violations of the safe region are detected by
\citep[Exogenous~Verification]{lasser2026exo}'s SPRT machinery within a
tail-bounded high-probability detection quantile $\Tbeta$ before irreversible
cascade (Lemmas~\ref{lem:5b}, \ref{lem:5e}, \ref{lem:5d}, \ref{lem:6}).  The
lead-time guarantee uses Wald's identity~\citep{wald1947sequential} and
Hoeffding--Azuma~\citep{hoeffding1963probability,azuma1967weighted} on the
clipped SPRT increment to derive $\Pr[T_{\mathrm{detect}} > \Tbeta] \leq
\beta$, with cascade-clock comparability via (C11.CLK) giving total Layer~2
failure $\leq \beta + \bclk + \bcal$.

\textbf{Layer 3 (acknowledged residuals).} Five named residuals are outside the
guarantee: (R1) channel-orthogonal restructuring --- strategies achieving
$\Delta r_K \geq r_{\mathrm{ext}}$ without producing any monitored four-channel
deviation; (R2) environment-witness-orthogonal manipulation --- environment
manipulations targeting only unmonitored exogenous variables and producing no
threshold-exceeding monitored shift; (R3) redundancy-dominated regime ---
deployments where redundant capabilities dominate cross-substrate cooperatives;
(R4) capability-targeting and coalition-internal shocks --- adversarial event
classes outside the substrate-targeting class; and (R5) calibration-exceeded
gap-growth --- adversarial strategies whose per-step gap growth exceeds the
calibrated $\rhogap$ from (C11).  We name these residuals explicitly rather
than gloss them.

The Goodhart slack between proxy and operational truth is bounded by an
intensive constant in Layer 1 and by an SPRT-tail-bounded intensive constant in
Layer 2. In neither case does the bound grow with the system's absolute
capability magnitude.

\subsection{Substrate identification and cooperative anchoring}
\label{sec:substrate_preview}

The composition's load-bearing piece is the canonical tripartite substrate
identification: \textbf{Human + AI + Formal-Operational}. These three are
jointly failure-correlation independent in the sense
\citep[Exogenous~Verification]{lasser2026exo}'s verification asymmetry argument
requires: prompt injection breaks neither human judgment nor cryptographic
commitments; cognitive bias breaks neither LLM inference nor ledger integrity;
trusted-setup failure breaks neither human reasoning nor LLM behavior.

A common alternative identification (``two LLMs and a human'') fails the test,
since LLMs share substantial adversarial-event surface through training-data
overlap, prompt-injection vulnerabilities, and architectural similarity. The
canonical tripartite identification makes the formal-operational layer (the
verification infrastructure itself) explicit as the third substrate, rather
than treating it as an undifferentiated ``safety machinery.''

This identification has a structural consequence
(\S\ref{sec:substrate_anchoring}): cross-substrate cooperative capabilities
involving the formal-operational layer are jointly produced and not
unilaterally replaceable. An agent cannot substitute a captured or weaker
verification layer because the cooperative's value comes from the human
substrate's actual participation in the specific shared infrastructure. This
\emph{cooperative-anchoring} property defeats the literal
replacement-verification attack on the deployment claim, narrowing the agent's
capture options to attacks \citep[Exogenous~Verification]{lasser2026exo}'s
existing defenses already address.

\subsection{Scope, residuals, and what we do not claim}
\label{sec:scope}

This paper's deployment claim is conditional. We are explicit about what each
conditional carries:

\textbf{The eleven operational invariants} ($I_1$--$I_{11}$,
\S\ref{sec:invariants}) are the deployment-time conditions. They include the
inherited \citep[Phase~Redundancy]{lasser2026phase} invariants (cross-substrate
redundancy $\rhomincross > \rho^*$, subsumption frequency $\rsub < r^*$,
Lyapunov error $L < \epsilon_{\mathrm{safe}}$), inherited
\citep[Revealed~Sacrifice]{lasser2026paper8a},~\citep[Need~Sufficiency]{lasser2026paper8b}
invariants (B-to-C ratio, HHI), and new operational machinery ($I_8$
environment-side substrate-exclusive witnesses, $I_9$ substrate-exclusivity
observability, $I_{10}$ coverage/materiality, $I_{11}$ latency bounds). Each
invariant is ledger-observable under
\citep[Exogenous~Verification]{lasser2026exo}'s verification infrastructure.

\textbf{The substrate-targeting scope.} The static safe region covers
substrate-targeting adversarial events specifically. Capability-targeting
attacks, coalition-internal corruption, and environment manipulation outside
the witness coverage are handled by separate machinery (the detection layer for
some, acknowledged residuals for others).

\textbf{The cooperative-overlap regime.} Lemma 5c's single-shock safe-region
result is established for deployments where the substrate-cooperative
structure is dominated by cross-substrate cooperatives (the canonical
tripartite case) rather than redundant capabilities. Sequential
counterfactual derivation for the redundancy regime is named as Lemma
5c-prime open work.

\textbf{The weaker inner-alignment condition.} The cooperative-anchoring
property requires the agent's effective objective to realize \emph{causally
grounded cooperative-outcome value}, not merely reward-visible cooperation
signals. This is more plausible than a strong substrate-aware mesa-objective
requirement~\citep{hubinger2019risks} but is not delivered by vanilla RLHF;
achieving it requires training-time discipline (delayed feedback, adversarial
fake-verification
examples, process supervision tied to real attestations).

\textbf{The Concentration-Gap structural conditions.} Theorem~1's Layer~1
binding via invariant $I_5$ requires four operationally auditable structural
conditions: representativeness (REP), bounded dispersion (DISP), coalition
closure (COAL), and Lipschitz embedding compatibility (EMB), collected as
Assumption~\ref{ass:concgap_conditions}. The conjunction of these four
conditions plays the role previously assigned to a monolithic
``HHI surrogate adequacy'' assumption; the companion
paper~\citep{lasser2026foundations} derives the Goodhart-slack bound under
these structural conditions, replacing what was previously an inherited
empirical-adequacy claim with four explicit auditable conditions plus the
HHI threshold itself. \S\ref{sec:conjecture_op} summarizes the structural
content; \citep[\S 4]{lasser2026foundations} contains the formal derivation.

\textbf{What we explicitly do not claim.}  This paper does not prove that
alignment pressure is universally reversed under substrate identification.  It
does not solve the inner-alignment problem: condition (C4) is assumed, and
requires additional training infrastructure to be bounded.  It does not
establish basin entry: whether a deployment can reach the cooperative-anchoring
attractor from arbitrary initial conditions requires empirical training-time
discipline from AI designers.  It does not bound the residual classes named in
Layer~3. \S\ref{sec:conditional_structure} works through the distinction
between conjectural conditions (argued, not proved), operational conditions
(verifiable from deployment state), and explicit scope restrictions.

\subsection{Paper roadmap}
\label{sec:roadmap}

\S\ref{sec:composition_setup} sets up the composition: how the theoretical
foundation stacks into a unified framework, with notation reconciliation across
the source papers. \S\ref{sec:invariants} defines the eleven operational
invariants with their threshold semantics and ledger-observable measurement
procedures. \S\ref{sec:lemmas} states and proves the ten compositional lemmas
(1--4, the five-part Lemma 5 family, and Lemma 6).
\S\ref{sec:main_theorem} states the three-layer deployment-safety theorem and
gives its proof.
\S\ref{sec:substrate_anchoring} develops the substrate identification and
cooperative-anchoring property in detail, including the three additional
invariants $I_9$--$I_{11}$ that bound non-substrate-targeting evasions.
\S\ref{sec:conjecture_op} summarizes the operationalization of
\citep[Microfoundation]{lasser2026micro}'s Concentration-Gap~Conjecture (the
Herfindahl--Hirschman Index as surrogate for optimization pressure) and points
to the companion paper~\citep{lasser2026foundations} for the formal scoped
discharge under the structural conditions of
Assumption~\ref{ass:concgap_conditions}.
\S\ref{sec:deployment_tooling} specifies what a deployer must instrument to
invoke the theorem. \S\ref{sec:worked_scenarios} walks through worked
deployment scenarios, including a clean case, per-invariant violations, an
adversarial-coalition scenario, and the canonical tripartite audit.
\S\ref{sec:discussion} discusses what this paper establishes and what it
defers, with explicit open questions for follow-up work.
